% !TeX root = ../main.tex
\clearpage
\section*{附录4\quad 数据来源、实际规模与关键参数}
\phantomsection\label{appendix:data}
\addcontentsline{toc}{section}{附录4：数据来源、实际规模与关键参数}

本附录按“实际进入结果的输入、只作背景或方法参考的资料、未准入资料”登记数据。规模按城市、建筑、时间窗口和目录工况分别计算；同一天气的原表、派生表和重复求解不叠加为独立样本。公开天气不称当地传感器实测，情景负荷不称计量热负荷，公共热原型与四楼合成参数不称现场校准。下列D01--D15与随文复核清单中的来源编号一致。

\subsection*{4.1\quad 实际使用的数据与规模}

原三城算例采用11664个连续城市小时输入行；多年风险层保留2020--2024年15组年度气象序列。外部计量确认另取六栋2017年建筑，设备专项另取两型号165对目录工况。三类材料分别用于检查连续服务、气象扰动、数据压缩和设备转换关系，不合并成一个现场实测数据集。

\begingroup
\setlength{\tabcolsep}{4pt}
\renewcommand{\arraystretch}{1.16}
\AtBeginEnvironment{longtable}{\fontsize{10.95}{13}\selectfont}
\begin{longtable}{@{}L{0.15\textwidth}L{0.24\textwidth}L{0.28\textwidth}L{0.25\textwidth}@{}}
\caption{实际输入对象、规模与结果用途}\label{tab:data-scale-r8p}\\
\toprule
数据类别 & 确用规模 & 来源及处理 & 结果用途与边界 \\
\midrule
\endfirsthead
\multicolumn{4}{l}{\small 表\ref{tab:data-scale-r8p}（续）}\\
\toprule
数据类别 & 确用规模 & 来源及处理 & 结果用途与边界 \\
\midrule
\endhead
\midrule
\multicolumn{4}{r}{\small 续下页}\\
\endfoot
\bottomrule
\endlastfoot
D01 原三城全季 & 乌鲁木齐4392 h；固原4392 h；武汉2880 h。合计11664行。 & NASA逐时天气＋18\,${}^{\circ}$C基温度时缩放的情景负荷；源容量、价格与资源曲线按配置设定。 & 三场景规划、锚定投资全季校核及相关策略对照；负荷不是当地热量表读数。 \\
D02 多年风险层 & 3城×5年＝15组；每城43848 h，合计131544个城市小时。风险权重12案、容量就绪检查60案年。 & 2020、2024各8784 h，其他年各8760 h；历史年度气象链标注LST，负荷由情景规则派生。 & 既有CVaR扫描与5气象年容量边界检查；60次检查不增加独立数据量。 \\
D03 六楼计量确认 & 六栋；52560个日历建筑小时；有效计量51955 h，联合天气有效51589 h；完整日46104 h。 & BDG2作者清洗版热水、蒸汽、冷水各两栋；保留缺失，按完整24 h日筛选。 & 1921个完整建筑日压缩为54个代表日（1296 h）；用于权重守恒、峰值/爬坡与W1检验。 \\
D04 固定热泵专项 & TCAV035BHE（83对），TCAV065BHE（82对），共165对$Q,P$；3城各2880 h，固定整机LP六案。 & TICA标识目录的第三方分发副本；20份公共原型的合成尺度为960 m$^2$。 & 以实际目录轴约束制热量和输入电功率；同服务范围对照能源费，不含采购和投资最优。 \\
D05--D06 公共原型与双楼 & 公共原型48 m$^2$；双楼四案，每案2楼×2880 h，同一武汉天气。 & BESTEST600规定几何及派生1R1C参数；B楼热容倍率和共享配电上限为研究设定。 & 分开检查热容轴、配电轴与每楼状态；四案没有新增四份独立天气或现场建筑。 \\
D07 四楼热状态 & 四栋合成建筑共44000 m$^2$；服务年8784 h，冬季2928 h；分别35136、11712个建筑状态小时。 & 武汉公开天气＋分别登记的面积、H、C、占用和舒适边界。 & C4四楼逐时状态和建筑、站点、源端账本；冬季位于服务年内，不加总为独立样本。 \\
D08 参考管网 & 72 h；近、远两支路，每支5个有限稳态模式；软惩罚关闭/开启两案。 & 取三城设备共同轴（2024-11-16起，见表\ref{tab:data-window-r8p}）的前72 h；几何、水物性、限值与模式均为设计参数。 & 末端温位、压降和源热--管损--送热--泵电核对；未作实测管网或冷启动验证。 \\
D09--D10 费用与排放 & 2024年11月至2025年3月，5个月×4时段＝20个完整电价；电力平均因子按各计算轴分别登记。 & 湖北公布能源价在账户适用条件下使用；排放为2023年平均因子参考。 & 专项共用湖北价格/因子作为方法轴；原三城仍保留自己的历史情景价格与国家平均因子。 \\
\end{longtable}

\subsection*{4.2\quad 时间窗口、单位与计量筛选}

所有逐时状态更新的时间步为1 h。三城原窗口按北京时间小时标签读取，起点含、终点不含；如下表列出的终点均为排他端点。固原窗口是论文情景运行边界，不能以4392 h代替当地行政供暖期。武汉窗口借用了2024年武昌及周边项目服务公告，也不是武汉全市法定供暖期。
\begin{longtable}{@{}L{0.2\textwidth}L{0.19\textwidth}L{0.19\textwidth}L{0.07\textwidth}L{0.27\textwidth}@{}}
\caption{连续运行窗口及计数规则}\label{tab:data-window-r8p}\\
\toprule
对象 & 开始（含） & 结束（不含） & 小时 & 窗口身份 \\
\midrule
\endfirsthead
\multicolumn{5}{l}{\small 表\ref{tab:data-window-r8p}（续）}\\
\toprule
对象 & 开始（含） & 结束（不含） & 小时 & 窗口身份 \\
\midrule
\endhead
\midrule
\multicolumn{5}{r}{\small 续下页}\\
\endfoot
\bottomrule
\endlastfoot
乌鲁木齐原全季 & 2024-10-10 00:00 & 2025-04-11 00:00 & 4392 & 历史情景输入 \\
固原原全季 & 2024-10-15 00:00 & 2025-04-16 00:00 & 4392 & 历史情景输入；不等同11月1日起公告 \\
武汉原全季／设备共同轴 & 2024-11-16 00:00 & 2025-03-16 00:00 & 2880 & 设备三城专项共同采用这一个窗口 \\
C4四楼服务年 & 2023-07-01 00:00 & 2024-07-01 00:00 & 8784 & 跨闰日服务年；不等于2023日历年 \\
C4四楼冬季 & 2023-11-15 00:00 & 2024-03-16 00:00 & 2928 & 服务年内子窗口 \\
72 h参考管网 & 2024-11-16 00:00 & 2024-11-19 00:00 & 72 & 设备共同轴的前三日 \\
\end{longtable}

NASA请求采用城市参考点：乌鲁木齐$(43.8256^{\circ}\mathrm N,87.6168^{\circ}\mathrm E)$、固原$(36.0159^{\circ}\mathrm N,106.2426^{\circ}\mathrm E)$、武汉$(30.5928^{\circ}\mathrm N,114.3055^{\circ}\mathrm E)$。温度为$^{\circ}$C，相对湿度原字段为\%，风速为m/s。太阳辐照按存档字段单位处理：旧年度表标为kWh/(m$^2\cdot$h)，后续UTC表标为Wh/m$^2$的小时累计量；不能只改字段名而跳过单位核对。NASA温度等气象与辐射使用不同网格产品，不能把一个统一空间分辨率转写为全部字段或现场精度。

多年风险结果对应原LST气象链；2026年10月保存的UTC重建包属于后续核验材料。C4使用武汉2023、2024年旧天气，并做了专项时间转换核对。时钟、单位与文件身份核对保证读入一致，不能据此声称气象精度或本地热负荷已完成校准。

六栋计量建筑均未参与此前开发样本；site仍与开发数据有重合。使用作者已规范化的小时kWh字段，除以1 h得到小时平均kW，不再按原始燃料或MJ重做一次单位换算。计量与联合温度覆盖率均须达到95\%；非负、有限且24 h完整的日才参与压缩，其余有效小时仍保存在损失核对中。热水通道的采暖与生活热水分界并未逐楼查明。
\begin{longtable}{@{}L{0.07\textwidth}L{0.36\textwidth}L{0.115\textwidth}L{0.115\textwidth}L{0.105\textwidth}L{0.1\textwidth}@{}}
\caption{六栋2017年计量样本的实际筛选数量}\label{tab:data-bdg2-r8p}\\
\toprule
通道 & 建筑ID & 有效计量h & 联合天气h & 完整日$N$ & 代表日$M$ \\
\midrule
\endfirsthead
\multicolumn{6}{l}{\small 表\ref{tab:data-bdg2-r8p}（续）}\\
\toprule
通道 & 建筑ID & 有效计量h & 联合天气h & 完整日$N$ & 代表日$M$ \\
\midrule
\endhead
\midrule
\multicolumn{6}{r}{\small 续下页}\\
\endfoot
\bottomrule
\endlastfoot
热水 & \texttt{Bobcat\_education\_Alissa} & 8641 & 8617 & 320 & 8 \\
热水 & \texttt{Eagle\_assembly\_Ian} & 8638 & 8636 & 345 & 10 \\
蒸汽 & \texttt{Bull\_assembly\_Brandon} & 8566 & 8531 & 304 & 9 \\
蒸汽 & \texttt{Cockatoo\_assembly\_Genoveva} & 8647 & 8405 & 276 & 9 \\
冷水 & \texttt{Bobcat\_education\_Angela} & 8721 & 8697 & 338 & 9 \\
冷水 & \texttt{Bull\_assembly\_Dorethea} & 8742 & 8703 & 338 & 9 \\
\end{longtable}

K=6是基础聚类请求数，补入极端日后实际M为8--10。六楼合计1921个完整日，保留54个代表日；不同权重方法使用相同已选24 h曲线。硬权重候选以总权重天数与保留域能量作等式约束，独立复核残差，再比较峰值、爬坡及负荷持续曲线的W1。研究统计门槛下的4/6准入只支持该压缩检验，不能转作工程调度、费用、碳或室温认证。

\subsection*{4.3\quad 关键输入参数与来源关系}

表\ref{tab:data-parameters-r8p}给出影响供热量、状态、费用和工况边界的主要数值。原情景参数仍按原配置登记；有公开来源的原值和派生值分别说明。没有追到项目测量或单项报价的参数保留情景身份，不给它补上“规范默认”或“厂家实测”的名义。
\begin{longtable}{@{}L{0.16\textwidth}L{0.32\textwidth}L{0.23\textwidth}L{0.21\textwidth}@{}}
\caption{主要数值输入、计算关系与资料身份}\label{tab:data-parameters-r8p}\\
\toprule
参数组 & 实际登记数值 & 使用方式 & 来源与范围 \\
\midrule
\endfirsthead
\multicolumn{4}{l}{\small 表\ref{tab:data-parameters-r8p}（续）}\\
\toprule
参数组 & 实际登记数值 & 使用方式 & 来源与范围 \\
\midrule
\endhead
\midrule
\multicolumn{4}{r}{\small 续下页}\\
\endfoot
\bottomrule
\endlastfoot
原三城温位 & 乌鲁木齐、固原75/45\,${}^{\circ}$C；武汉45\,${}^{\circ}$C。 & 供热品位与需求通道配置 & D01情景温位；不是已测管网供回水。 \\
原三城设计室外温度 & 乌鲁木齐$-26$、固原$-26$、武汉$-7$\,${}^{\circ}$C。 & 既有设计/保供校核输入 & D01历史情景值；不冒充D11规范表。 \\
原三城财务轴 & 折现率5\%，资产寿命20年；时间步1 h。 & 容量费用年化、运行量累计 & D01研究配置；不是实际融资合同。 \\
原三城碳价／储备 & 碳价0.10/0.10/0.12元/kg；保供容量储备25\%/15\%/12\%；关键峰段12/12/6 h。 & 按乌鲁木齐／固原／武汉顺序 & D01情景政策和保供轴；不称行政统一标准。 \\
原三城电价范围 & 0.246848--0.691；0.19173--0.35607；0.4953--1.1795元/kWh。 & 三个历史输入CSV的实际最小、最大值 & D01保留原情景价；与D09另表公布价不混用。 \\
旧源与接入费用 & 空气热泵1700--1850、地源3100--3150、电锅炉560--650元/kW；电接入新增容量费用820元/kW。 & 所列为旧三城配置中对应技术的资本参数 & D01情景报价轴，未建立当地项目逐项原始报价关系；不转为本轮新采购事实。 \\
公共原型H、C & 48 m$^2$，129.6 m$^3$；H=0.1043 kW/K，C=0.535576 kWh/K；内增益0.2 kW。 & BESTEST600规定几何，围护/渗透项与室内侧材料归并 & D05公开原型派生；不是地方楼体辨识。 \\
热泵合成尺度 & 20份原型，960 m$^2$；H=2.086 kW/K，C=10.71152 kWh/K。 & 对原型H、C线性复制为研究尺度 & D04/D05；不称20户实测或已识别等效楼。 \\
热泵服务与资产边界 & 室温20--27\,${}^{\circ}$C，初末20\,${}^{\circ}$C；出水45\,${}^{\circ}$C；1整模块；电热90 kW、独立制冷80 kW、配电100 kW。 & 固定设备六案共享服务包络；独立冷机EER=3，电热效率1 & D04固定理想资产；目录没有提供实际部分负荷/启停曲线。 \\
两型号工况 & 035型号83对，065型号82对$Q,P$；目录环境轴$-26$至55\,${}^{\circ}$C。 & 逐点核对；$Q,P$分别插值，再计算COP=$Q/P$ & D04打印轴内且所需角点齐全才准入，空白角点和外推拒绝；湿度/化霜信息未补齐。 \\
双楼2×2参数 & A：H=0.1043，C=0.535576；B：H=0.15645，C=0.535576或2.142304。H单位kW/K，C单位kWh/K。 & B楼H保持1.5倍，C设1倍或4倍；共享电上限另设开/关 & D06设计轴；没有更换实际墙体或现场测量。 \\
双楼公用电上限 & 6.4796375 kW；按$2.5H_A[20-\min(T_{\rm out})]$预先确定。 & 四案同2880 h，初末20\,${}^{\circ}$C、服务20--27\,${}^{\circ}$C & D06；原理想热/冷设备各1000 kW、效率/EER均1，共享配电是唯一联系。 \\
参考管网几何 & 近支路单程30 m、远支路150 m；内径0.015 m；粗糙度$5\times10^{-5}$ m。 & 供回双腿热损分别计算 & D08设计值，不代表某现场DN管径和竣工长度。 \\
参考管网热边界 & 源供65、地温5、末端供温下限55、回温上限45\,${}^{\circ}$C；每腿散热系数0.3 W/(m$\cdot$K)。 & 两支路同一硬边界；送热与热损分别计量 & D08设计限值；不能宣称规范自动给定。 \\
参考管网水物性／泵 & 密度985 kg/m$^3$，比热4.18 kJ/(kg$\cdot$K)，黏度0.00055 Pa$\cdot$s；泵效率0.72。 & 两支路独立泵，局部阻力系数2。 & D08固定设计参考；不是随温实测物性。 \\
参考管网模式／容量 & 各支路质量流量0.03、0.06、0.08、0.12、0.18 kg/s；模式回温40\,${}^{\circ}$C；源40 kW，配电50 kW。 & 每支每小时最多一个稳态开态，另有占空边界 & D08有限候选模式；不作连续全水力可行域认证。 \\
参考管网硬／软限值 & 压降300 Pa/m；流速3 m/s；单支泵0.25 kW、瞬时热输入30 kW；温差目标20\,${}^{\circ}$C、软带1\,${}^{\circ}$C。 & 硬限筛除模式；软惩罚0或0.2元/(\,${}^{\circ}$C$\cdot$开态h) & D08设计机制测试；软惩罚与源/泵能源费单列。 \\
电力平均排放轴 & 原三城0.5306；公共原型等专项共用湖北0.4044 kgCO$_2$/kWh。 & 2023年国家／湖北年度平均因子 & D10；不是2024/25逐时边际排放。新疆0.6021、宁夏0.6187为待重跑参考，未替换旧结果。 \\
\end{longtable}

H由已登记围护传热、渗透等项求和；C由选定室内侧材料的$A d\rho c_p$归并并换算为kWh/K。它们的量纲与1R1C方程一致，但源于规定原型与归并假设，未用三城实楼室温及实际供热输入作辨识。四楼C4与公共原型双楼是两组独立设计，不能混称同一批建筑。
\begin{longtable}{@{}L{0.22\textwidth}L{0.13\textwidth}L{0.13\textwidth}L{0.15\textwidth}L{0.13\textwidth}L{0.13\textwidth}@{}}
\caption{C4四栋合成建筑的面积和热参数}\label{tab:data-c4-parameters-r8p}\\
\toprule
建筑 & 面积m$^2$ & H kW/K & C kWh/K & 热损W/(m$^2$K) & 热容Wh/(m$^2$K) \\
\midrule
\endfirsthead
\multicolumn{6}{l}{\small 表\ref{tab:data-c4-parameters-r8p}（续）}\\
\toprule
建筑 & 面积m$^2$ & H kW/K & C kWh/K & 热损W/(m$^2$K) & 热容Wh/(m$^2$K) \\
\midrule
\endhead
\midrule
\multicolumn{6}{r}{\small 续下页}\\
\endfoot
\bottomrule
\endlastfoot
\texttt{apartments\_a} & 18000 & 16.2 & 1260 & 0.9 & 70 \\
\texttt{office\_a} & 7000 & 7.7 & 315 & 1.1 & 45 \\
\texttt{apartments\_b} & 14000 & 14 & 980 & 1.0 & 70 \\
\texttt{retail\_b} & 5000 & 6 & 175 & 1.2 & 35 \\
合计 & 44000 & 43.9 & 2730 & -- & -- \\
\end{longtable}

上表全部标记为\texttt{synthetic\_assumption\_not\_calibrated}。初温20\,${}^{\circ}$C，基础上下界18--28\,${}^{\circ}$C；实际计算同时使用每楼逐时占用、增益及温度边界列。面积、H、C与舒适规则提供可复算的异质设计对象，不能作为四栋当地建筑的实测档案。

\subsection*{4.4\quad 条件电价与费用边界}

专项实验选择湖北普通工商业代理购电、单一制、不满1千伏的公布能源价行。适用账户需满足已登记供电容量区间与至少12个月的选择期限等条件，论文未核实任一具体账户的资格或接入。公布的完整能源价直接代入小时电量；不再乘一次峰谷系数。配电模型上限不是账户登记容量，设备380 V标签也不能代替计费供电电压身份。
\begin{longtable}{@{}L{0.11\textwidth}L{0.13\textwidth}L{0.13\textwidth}L{0.13\textwidth}L{0.13\textwidth}L{0.27\textwidth}@{}}
\caption{条件参考轴的公布能源电价（元/kWh）}\label{tab:data-tariff-r8p}\\
\toprule
月份 & 尖峰 & 高峰 & 平段 & 低谷 & 原表定位 \\
\midrule
\endfirsthead
\multicolumn{6}{l}{\small 表\ref{tab:data-tariff-r8p}（续）}\\
\toprule
月份 & 尖峰 & 高峰 & 平段 & 低谷 & 原表定位 \\
\midrule
\endhead
\midrule
\multicolumn{6}{r}{\small 续下页}\\
\endfoot
\bottomrule
\endlastfoot
2024-11 & 1.1724 & 1.0188 & 0.7762 & 0.5187 & 11月通知第1页 \\
2024-12 & 1.2908 & 1.0368 & 0.7927 & 0.5187 & 12月通知第1页 \\
2025-01 & 1.1995 & 0.9744 & 0.7581 & 0.5153 & 2025年合订通知第12页 \\
2025-02 & 1.0553 & 0.9313 & 0.7353 & 0.5273 & 2025年合订通知第11页 \\
2025-03 & 1.1144 & 0.9726 & 0.7484 & 0.5105 & 2025年合订通知第10页 \\
\end{longtable}

本窗口小时归属为：尖峰18、19时；高峰16、17、20--23时；平段6--11、14、15时；低谷0--5、12、13时，各时间戳标记随后1 h。电价含公布列中的线损、输配、基金及系统运行等能源价组成；实际功率因数调整、其他服务、结算修正、采购、维护与配电投资不在专项能源费中。热泵三城对照共用该条件轴以隔离天气/设备差异，不预测乌鲁木齐、固原的本地账单。

\subsection*{4.5\quad 原始来源与仅参考资料}

来源链接指向数据或公布材料本身，详细文件路径、物理页码、原响应、筛选规则与SHA-256保存在\path{review/data_inventory_r8p.json}。指纹只核对保存身份，不认证资料准确性、现场适用性或方案最优。下面将实际数值来源与背景参考并列，避免以下载库存扩大论文使用量。
\begin{longtable}{@{}L{0.17\textwidth}L{0.44\textwidth}L{0.31\textwidth}@{}}
\caption{来源链接、确用内容与参考边界}\label{tab:data-source-r8p}\\
\toprule
来源 & 材料与定位 & 实际用途／未采用范围 \\
\midrule
\endfirsthead
\multicolumn{3}{l}{\small 表\ref{tab:data-source-r8p}（续）}\\
\toprule
来源 & 材料与定位 & 实际用途／未采用范围 \\
\midrule
\endhead
\midrule
\multicolumn{3}{r}{\small 续下页}\\
\endfoot
\bottomrule
\endlastfoot
D01--D02\newline NASA POWER & \href{https://power.larc.nasa.gov/docs/services/api/temporal/hourly/}{官方逐时数据说明}；各城市请求、原响应及UTC/LST身份随清单保存。 & 实际气象输入；温度、湿度、风、辐照；负荷与太阳热可用率是后续派生。 \\
D03\newline BDG2 & \href{https://doi.org/10.1038/s41597-020-00712-x}{Miller等，Scientific Data 7:368（2020）}；\href{https://github.com/buds-lab/building-data-genome-project-2/tree/9b97ccbe90096aff42ed4fd6493bf7ae692d7118}{作者固定版本数据仓库}。 & 确用六栋2017年计量通道；原型建筑H/C、室温、设备COP与当地账单不由此获取。 \\
D04\newline TICA目录 & \href{https://dec.com.pk/assets/images/products/tica/3-Inverter-modular-chiller.pdf}{24页产品目录分发副本}，物理第15页$Q,P$网格；第14页额定工况另记。 & 两型号83/82对打印值经逐点核对；原刊端该具体表版本尚未认证。 \\
D05--D06\newline BESTEST & \href{https://www.iea-ebc.org/Data/publications/EBC_Annex_21_BESTEST.pdf}{1995年NREL/TP-472-6231公共原型报告}，物理第34、36、39、41、48--50、53页。 & 规定几何和材料用于H/C派生；没有当地辨识或BESTEST软件认证。 \\
D09\newline 五个月电价 & \href{http://www.yingcheng.gov.cn/u/cms/yingcheng/202410/29154032ntk8.pdf}{2024年11月表}、\href{http://www.yingcheng.gov.cn/u/cms/yingcheng/202412/04105911f9b9.pdf}{12月表}及\href{http://www.hb.sgcc.com.cn/html/files/2026-03/06/20260306144020453109213.pdf}{2025年合订表}。 & 前两份为政府转载；合订表为国网湖北公布。选同一单一制、不满1kV行，物理页码见表。 \\
D09\newline 账户与时段规则 & \href{https://www.liangzh.gov.cn/lzhxxgk/xxgkml/ggqsydwxxgk/sdq/202508/P020250825377454860287.pdf}{鄂发改价管〔2024〕77号}及\href{https://www.liangzh.gov.cn/lzhxxgk/xxgkml/ggqsydwxxgk/sdq/202508/P020250825378015593092.pdf}{〔2023〕163号}原件政府转载。 & 77号第2页范围/时段、第3页生效日；163号第2页账户选择条件。按历史窗口作条件引用。 \\
D10\newline 电力因子 & \href{https://www.mee.gov.cn/xxgk2018/xxgk/xxgk01/202512/t20251231_1139517.html}{生态环境部、国家统计局2025年第47号公告}及2023年因子原附件。 & 实际国家0.5306／专项湖北0.4044；省级重跑建议表不是全部旧结果的已用输入。 \\
D11\newline 设计气候表 & \href{https://gf.cabr-fire.com/m/article-12078.htm}{GB 50736--2012附录A刊载原表}，1971--2000气候统计。 & 仅参考：供暖设计温度乌鲁木齐$-19.7$、固原$-13.2$、武汉$-0.3$\,${}^{\circ}$C；不替换原情景值，不宣布现行规范审查完成。 \\
D12\newline 城市规划 & \href{https://www.urumqi.gov.cn/wlmqs/c119242/202506/282e876f9e074a4f837ea68ae7b30f59.shtml}{乌鲁木齐2025--2027供热规划草案}；\href{https://www.nxgy.gov.cn/zwgk/zcwj/qtwj/jgzn_35170/202208/W020220822623460812528.pdf}{固原2021--2025清洁取暖规划}；\href{https://www.caidian.gov.cn/zwgk/xxgkml/zdxm/xmjys/202012/t20201201_1523964.shtml}{蔡甸2020年项目审批材料}。 & 场景背景参考；规划面积/容量不自动等于数值原型存量、竣工资产或本轮实绩。 \\
D13\newline OSM空间资料 & \href{https://www.openstreetmap.org/copyright}{OpenStreetMap来源与许可}；原CSV导出缺快照日期、查询范围和版本元数据。 & 三城建筑质心、道路等可作空间参考；现有审计未执行优化，不计入本论文已用管段或实楼H/C。 \\
D14\newline UTC重建包 & 保存三城15个年度、3个季节NASA请求和坐标索引，区别于旧年度LST输入。 & 用于来源与时钟核验；原返回含边界补齐，按最终窗口裁剪；不能据库存数量宣称已新跑年度实验。 \\
D15\newline 未准入与待补 & 39页营销资料、跨型号采购价格、没有原始凭证的口述性能；当地实楼测温、实际热输入、资产报价和连接审批。 & 不补作实测参数。目录空点、外推点、无适用账户或不明来源项保留拒绝/待核状态。 \\
\end{longtable}

\endgroup

完整设备存量、费用、资源曲线与候选管段见结构化清单。公开计量、目录工况和热原型分别检验压缩误差、设备转换与状态守恒；情景参数保留原身份，工程实施前另核现场凭证。
